\vspace{0.5em}\noindent\textbf{WHY DOES AWS STILL CHARGE AFTER STOPPING
EC2?}\par

A common misunderstanding when learning AWS is: \emph{EC2 is stopped =
AWS stops charging.}

However, through posts and messages in the group, I have seen many
people still receive bills even when their instance is in the
\emph{Stopped} state, or even after it has been \emph{Terminated}.

One real case was a friend of mine who accidentally left a NAT Gateway
running after doing a lab and received a bill of nearly \$250.

The issue is that EC2 is only one part of the system. The resources
created together with it do not always disappear automatically.

\vspace{0.5em}\noindent\textbf{Stopping EC2 only stops virtual machine compute
charges}\par

When an EC2 instance is in the Stopped state, AWS stops charging for the
compute part. However, the EBS disk is still kept so the machine can be
started again later, so the storage capacity is still charged.

A simple way to understand it:

\begin{itemize}
\tightlist
\item
  Stop = Turn off the machine but keep the disk.
\item
  Terminate = Delete the machine, but it is still necessary to check
  whether the EBS volume is configured to be deleted with the instance.
\end{itemize}

\vspace{0.5em}\noindent\textbf{NAT Gateway still charges even when EC2 is
stopped}\par

NAT Gateway is a separate resource in a VPC, so it does not stop
automatically with EC2.

AWS charges based on both:

\begin{itemize}
\tightlist
\item
  The number of hours the NAT Gateway exists.
\item
  The amount of data passing through it.
\end{itemize}

For example, in AWS pricing for some regions, a NAT Gateway may cost
around \$0.045 per hour (reference cost), equivalent to more than \$32
per month, even with almost no traffic.

If data passes through it, the cost increases further.

\vspace{0.5em}\noindent\textbf{Elastic IP is not free
either}\par

If an Elastic IP has been allocated but not released, that address may
continue to generate cost.

AWS currently charges for public IPv4 addresses whether they are in use
or idle. One address may only cost a few dollars per month, but if many
addresses are forgotten across multiple regions, the cost can add up
significantly.

\vspace{0.5em}\noindent\textbf{AWS Budget is not a hard spending
limit}\par

By default, AWS Budget is mainly used to send alerts. It does not
automatically lock the entire account when spending reaches the
configured amount.

If the system needs to respond automatically, Budget Actions must be
configured, for example:

\begin{itemize}
\tightlist
\item
  Stop a specific EC2 or RDS resource.
\item
  Attach an IAM policy to prevent creating more resources.
\item
  Send alerts through email or SNS.
\end{itemize}

\vspace{0.5em}\noindent\textbf{Reference checklist after each AWS
lab}\par

\begin{itemize}
\tightlist
\item
  Open Cost Explorer and group by Service and Usage Type.
\item
  Check EC2 \(\rightarrow\) Volumes and delete EBS volumes that are no
  longer needed.
\item
  Check VPC \(\rightarrow\) NAT Gateways.
\item
  Release Elastic IPs that are no longer used.
\item
  Check the correct region because resources in other regions can still
  generate cost.
\item
  Enable AWS Budget and Cost Anomaly Detection from the beginning.
\end{itemize}

The most important lesson I learned is: when deleting a service on AWS,
do not only ask: \emph{``Have I turned it off?''} Also ask:

\emph{``What supporting resources did this service create, and are they
still generating cost?''}

\texttt{\#AWS} \texttt{\#EC2} \texttt{\#CloudComputing}
\texttt{\#AWSCostOptimization}
